% book05.tex --- Book V of Euclid's Elements: Eudoxean Theory of Proportion.
%
% All 25 propositions encoded. The proofs are essentially algebraic
% manipulations of the equimultiples definition (V.5), so the dependency
% structure is dense within Book V itself rather than reaching back into
% Book I. Cross-book dependencies on V are heavy: Book VI (similarity)
% and Book XII (exhaustion) both rely on this machinery.
%
% Wording follows Heath (1908).

\section{Book V --- Eudoxean Theory of Proportion}
\label{sec:book-V}

\begin{claim}[Proposition V.1: Multiplication is distributive over magnitudes]
\label{prop:V.1}
If there be any number of magnitudes whatever which are, respectively,
equimultiples of any magnitudes equal in multitude, then, whatever
multiple one of the magnitudes is of one, that multiple also will all
be of all.
\end{claim}
\begin{evidence}[Proof of V.1]
\label{ev:V.1}
Each magnitude is a sum of $m$ copies of the corresponding base
magnitude.  Sum across the $n$ magnitudes; by Common Notion 2 the
total is $m$ copies of the sum.
\dependson{V.1}{cn:2}
\dependson{V.1}{def:V.2}
\end{evidence}

\begin{claim}[Proposition V.2: Equimultiples sum to equimultiples]
\label{prop:V.2}
If a first magnitude be the same multiple of a second that a third is
of a fourth, and a fifth also be the same multiple of the second that
a sixth is of the fourth, then the sum of the first and fifth will
also be the same multiple of the second that the sum of the third and
sixth is of the fourth.
\end{claim}
\begin{evidence}[Proof of V.2]
\label{ev:V.2}
If $a = mb$ and $c = md$, $e = nb$ and $f = nd$, then $a + e = (m+n)b$
and $c + f = (m+n)d$ by Common Notion 2.
\dependson{V.2}{V.1}
\dependson{V.2}{cn:2}
\dependson{V.2}{def:V.2}
\end{evidence}

\begin{claim}[Proposition V.3: Multiple of a multiple is a multiple]
\label{prop:V.3}
If a first magnitude be the same multiple of a second that a third is
of a fourth, and if equimultiples be taken of the first and third,
they will also be equimultiples respectively, the one of the second
and the other of the fourth.
\end{claim}
\begin{evidence}[Proof of V.3]
\label{ev:V.3}
If $a = mb$ and $c = md$, then $na = (nm)b$ and $nc = (nm)d$ by
iterated application of V.1.
\dependson{V.3}{V.1}
\dependson{V.3}{def:V.2}
\end{evidence}

\begin{claim}[Proposition V.4: Proportionality preserved under equimultiples]
\label{prop:V.4}
If a first magnitude have to a second the same ratio as a third to a
fourth, any equimultiples whatever of the first and third will also
have the same ratio to any equimultiples whatever of the second and
fourth respectively, taken in corresponding order.
\end{claim}
\begin{evidence}[Proof of V.4]
\label{ev:V.4}
By definition V.5, the original proportion gives an equimultiple
relation across all multipliers.  Substituting $m a$ for $a$ and $m c$
for $c$ throughout simply rescales the test multipliers; the test
itself still passes for the rescaled magnitudes.
\dependson{V.4}{V.3}
\dependson{V.4}{def:V.5}
\end{evidence}

\begin{claim}[Proposition V.5: Subtraction of equimultiples]
\label{prop:V.5}
If a magnitude be the same multiple of a magnitude that a subtracted
part is of a subtracted part, the remainder also will be the same
multiple of the remainder that the whole is of the whole.
\end{claim}
\begin{evidence}[Proof of V.5]
\label{ev:V.5}
If $a = mb$ and $a' = mb'$ with $a' < a$, then $a - a' = m(b - b')$
by Common Notion 3 applied to each of the $m$ copies.
\dependson{V.5}{cn:3}
\dependson{V.5}{def:V.2}
\end{evidence}

\begin{claim}[Proposition V.6: Difference of equimultiples is again an equimultiple]
\label{prop:V.6}
If two magnitudes be equimultiples of two magnitudes, and any
magnitudes subtracted from them be equimultiples of the same, the
remainders also are either equal to the same or equimultiples of them.
\end{claim}
\begin{evidence}[Proof of V.6]
\label{ev:V.6}
With $a = mb$, $c = md$, and subtractions $a' = nb$, $c' = nd$:
$a - a' = (m-n)b$ and $c - c' = (m-n)d$ by Common Notion 3.  If $m = n$
the remainders are zero (equal); otherwise both remainders are
$(m-n)$-fold of $b$, $d$ respectively.
\dependson{V.6}{V.5}
\dependson{V.6}{cn:3}
\end{evidence}

\begin{claim}[Proposition V.7: Equal magnitudes have the same ratio to a third]
\label{prop:V.7}
Equal magnitudes have to the same the same ratio, as also has the
same to equal magnitudes.
\end{claim}
\begin{evidence}[Proof of V.7]
\label{ev:V.7}
Let $a = b$ and $c$ arbitrary.  For any equimultiples $ma$, $mb$ of
$a$, $b$, and $nc$ of $c$, the test of V.5 succeeds vacuously because
$ma = mb$.  Therefore $a:c = b:c$ and $c:a = c:b$.
\dependson{V.7}{def:V.5}
\dependson{V.7}{cn:1}
\end{evidence}

\begin{claim}[Proposition V.8: Greater magnitude has greater ratio]
\label{prop:V.8}
Of unequal magnitudes the greater has to the same a greater ratio
than the less has, and the same has to the less a greater ratio than
it has to the greater.
\end{claim}
\begin{evidence}[Proof of V.8]
\label{ev:V.8}
Let $a > b$.  Pick a multiplier $n$ such that $n(a-b) > c$ (the
Archimedean property of magnitudes, Definition V.4) and an $m$ such
that $mc$ falls between $nb$ and $na$.  Then $na > mc$ but $nb < mc$;
by Definition V.7 this is precisely the assertion $a:c > b:c$.
\dependson{V.8}{def:V.4}
\dependson{V.8}{def:V.7}
\end{evidence}

\begin{claim}[Proposition V.9: Magnitudes with the same ratio to a third are equal]
\label{prop:V.9}
Magnitudes which have the same ratio to the same are equal to one
another; and magnitudes to which the same has the same ratio are
equal.
\end{claim}
\begin{evidence}[Proof of V.9]
\label{ev:V.9}
Contrapositive of V.8: if $a \neq b$, then $a:c \neq b:c$.  Hence if
$a:c = b:c$ then $a = b$.  Same argument with $c$ as antecedent.
\dependson{V.9}{V.8}
\end{evidence}

\begin{claim}[Proposition V.10: Greater ratio implies greater antecedent]
\label{prop:V.10}
Of magnitudes which have a ratio to the same, that which has a greater
ratio is greater; and that to which the same has a greater ratio is
less.
\end{claim}
\begin{evidence}[Proof of V.10]
\label{ev:V.10}
Same contrapositive of V.8: if $a:c > b:c$ then $a > b$, since if
$a \le b$ then $a:c \le b:c$ by V.7 or V.8 applied in reverse.
\dependson{V.10}{V.7}
\dependson{V.10}{V.8}
\end{evidence}

\begin{claim}[Proposition V.11: Transitivity of ratios]
\label{prop:V.11}
Ratios which are the same with the same ratio are also the same with
one another.
\end{claim}
\begin{evidence}[Proof of V.11]
\label{ev:V.11}
If $a:b = c:d$ and $c:d = e:f$, then for any equimultiples the same
inequality test holds for $(a,b)$ as for $(c,d)$, and that same test
holds for $(c,d)$ as for $(e,f)$; hence by transitivity of the
inequality test, the test holds for $(a,b)$ versus $(e,f)$.
\dependson{V.11}{def:V.5}
\end{evidence}

\begin{claim}[Proposition V.12: Sum of antecedents to sum of consequents]
\label{prop:V.12}
If any number of magnitudes be proportional, as one of the antecedents
is to one of the consequents, so will all the antecedents be to all
the consequents.
\end{claim}
\begin{evidence}[Proof of V.12]
\label{ev:V.12}
Let $a_i : b_i$ all equal $r$ in the sense of Definition V.5.  For
any test multipliers $m$, $n$ the sign of $m a_i - n b_i$ is the same
for every $i$; therefore the sign of $m \sum a_i - n \sum b_i$ is the
same too.  By V.5 this is the equimultiples test for $\sum a_i : \sum
b_i = r$.
\dependson{V.12}{V.1}
\dependson{V.12}{V.2}
\dependson{V.12}{def:V.5}
\end{evidence}

\begin{claim}[Proposition V.13: Substitution into a greater ratio]
\label{prop:V.13}
If a first magnitude have to a second the same ratio as a third to a
fourth, and the third have to the fourth a greater ratio than a fifth
has to a sixth, the first will also have to the second a greater ratio
than the fifth to the sixth.
\end{claim}
\begin{evidence}[Proof of V.13]
\label{ev:V.13}
Combine V.11 (sameness transitivity) with Definition V.7 (greater
ratio): the witness equimultiples for $c:d > e:f$ work for $a:b$ via
the V.5 sameness of $a:b$ and $c:d$.
\dependson{V.13}{V.11}
\dependson{V.13}{def:V.5}
\dependson{V.13}{def:V.7}
\end{evidence}

\begin{claim}[Proposition V.14: Ordering of magnitudes follows ordering of ratios]
\label{prop:V.14}
If a first magnitude have to a second the same ratio as a third to a
fourth, and the first be greater than the third, the second will also
be greater than the fourth; and if equal, equal; and if less, less.
\end{claim}
\begin{evidence}[Proof of V.14]
\label{ev:V.14}
By V.8 and V.13: if $a > c$, then $a:b > c:b$.  Combined with $a:b =
c:d$ (the hypothesis), V.13 gives $c:d > c:b$, whence $b > d$ by V.10.
\dependson{V.14}{V.8}
\dependson{V.14}{V.10}
\dependson{V.14}{V.13}
\end{evidence}

\begin{claim}[Proposition V.15: Parts have the same ratio as multiples]
\label{prop:V.15}
Parts have the same ratio as the same multiples of them taken in
corresponding order.
\end{claim}
\begin{evidence}[Proof of V.15]
\label{ev:V.15}
If $A = mB$ and $C = mD$, group the $m$ copies of $B$ and $D$ in
parallel.  Each parallel pair $(B_i, D_i)$ satisfies $B_i : D_i = B :
D$ (identity), and V.12 sums these to give $A : C = B : D$.
\dependson{V.15}{V.12}
\dependson{V.15}{def:V.2}
\end{evidence}

\begin{claim}[Proposition V.16: Alternation of proportionals]
\label{prop:V.16}
If four magnitudes be proportional, they will also be proportional
alternately.
\end{claim}
\begin{evidence}[Proof of V.16]
\label{ev:V.16}
Let $a : b = c : d$.  Test $a : c$ against $b : d$ with equimultiples
$ma$, $mc$, $nb$, $nd$: by V.4 the original proportion lifts to $ma :
mb = nc : nd$, and by V.15 to $ma : nc = mb : nd$.  The sign of $ma -
nc$ matches the sign of $mb - nd$ for all $m$, $n$, which by V.5 is
the alternated proportion $a : c = b : d$.
\dependson{V.16}{V.4}
\dependson{V.16}{V.11}
\dependson{V.16}{V.15}
\dependson{V.16}{def:V.5}
\end{evidence}

\begin{claim}[Proposition V.17: Separation of proportions]
\label{prop:V.17}
If magnitudes composed be proportional, they will also be proportional
separando.
\end{claim}
\begin{evidence}[Proof of V.17]
\label{ev:V.17}
If $(a + b) : b = (c + d) : d$, then subtracting the consequents from
the antecedents using V.5 / V.6 gives $a : b = c : d$.
\dependson{V.17}{V.5}
\dependson{V.17}{V.6}
\dependson{V.17}{def:V.5}
\end{evidence}

\begin{claim}[Proposition V.18: Composition of proportions]
\label{prop:V.18}
If magnitudes separated be proportional, they will also be
proportional componendo.
\end{claim}
\begin{evidence}[Proof of V.18]
\label{ev:V.18}
The converse of V.17: if $a : b = c : d$, then by V.2 plus V.4
combined, $(a + b) : b = (c + d) : d$.
\dependson{V.18}{V.2}
\dependson{V.18}{V.4}
\dependson{V.18}{V.17}
\end{evidence}

\begin{claim}[Proposition V.19: Subtraction of proportionals]
\label{prop:V.19}
If, as a whole is to a whole, so is a part subtracted to a part
subtracted, the remainder will also be to the remainder as whole to
whole.
\end{claim}
\begin{evidence}[Proof of V.19]
\label{ev:V.19}
If $a : c = a' : c'$ with $a' < a$, $c' < c$, apply V.17 (separation)
to obtain $(a - a') : a' = (c - c') : c'$, and V.11 / V.16 to
re-express as $(a - a') : (c - c') = a : c$.
\dependson{V.19}{V.16}
\dependson{V.19}{V.17}
\end{evidence}

\begin{claim}[Proposition V.20: Ex aequali for three magnitudes]
\label{prop:V.20}
If there be three magnitudes, and others equal to them in multitude,
which taken two and two are in the same ratio, and if ex aequali the
first be greater than the third, the fourth will also be greater than
the sixth; and if equal, equal; and if less, less.
\end{claim}
\begin{evidence}[Proof of V.20]
\label{ev:V.20}
With $a : b = d : e$ and $b : c = e : f$, the relative order of $a$
versus $c$ matches the relative order of $d$ versus $f$ by repeated
application of V.13 and V.14 across the three pairs.
\dependson{V.20}{V.13}
\dependson{V.20}{V.14}
\end{evidence}

\begin{claim}[Proposition V.21: Ex aequali in perturbed proportion]
\label{prop:V.21}
If there be three magnitudes, and others equal to them in multitude,
which taken two and two together are in the same ratio, and the
proportion of them be perturbed, then if ex aequali the first be
greater than the third, the fourth will also be greater than the
sixth.
\end{claim}
\begin{evidence}[Proof of V.21]
\label{ev:V.21}
A perturbed version of V.20: with $a : b = e : f$ and $b : c = d : e$
the ratio chain still imposes the same ordering between $a$ versus
$c$ and $d$ versus $f$, by V.13 and V.14.
\dependson{V.21}{V.13}
\dependson{V.21}{V.14}
\dependson{V.21}{def:V.18}
\end{evidence}

\begin{claim}[Proposition V.22: Transitivity of ex aequali]
\label{prop:V.22}
If there be any number of magnitudes whatever, and others equal to
them in multitude, which taken two and two together are in the same
ratio, they will also be in the same ratio ex aequali.
\end{claim}
\begin{evidence}[Proof of V.22]
\label{ev:V.22}
Apply V.20 inductively across the chain $a_1 : a_2 = b_1 : b_2$, $a_2
: a_3 = b_2 : b_3$, $\dots$, $a_{n-1} : a_n = b_{n-1} : b_n$ to
conclude $a_1 : a_n = b_1 : b_n$.
\dependson{V.22}{V.20}
\dependson{V.22}{def:V.17}
\end{evidence}

\begin{claim}[Proposition V.23: Ex aequali in perturbed proportion (general)]
\label{prop:V.23}
If there be any number of magnitudes whatever, and others equal to
them in multitude, which taken two and two together are in the same
ratio, and the proportion of them be perturbed, they will also be in
the same ratio ex aequali.
\end{claim}
\begin{evidence}[Proof of V.23]
\label{ev:V.23}
Inductive form of V.21.
\dependson{V.23}{V.21}
\dependson{V.23}{def:V.18}
\end{evidence}

\begin{claim}[Proposition V.24: Sums of antecedents are proportional]
\label{prop:V.24}
If a first magnitude have to a second the same ratio as a third has
to a fourth, and also a fifth have to the second the same ratio as a
sixth to the fourth, the first and fifth added together will have to
the second the same ratio as the third and sixth have to the fourth.
\end{claim}
\begin{evidence}[Proof of V.24]
\label{ev:V.24}
With $a : b = c : d$ and $e : b = f : d$: by V.12 applied to the two
proportions side-by-side, $(a + e) : b = (c + f) : d$.
\dependson{V.24}{V.12}
\dependson{V.24}{V.22}
\end{evidence}

\begin{claim}[Proposition V.25: Sum of extremes exceeds sum of means]
\label{prop:V.25}
If four magnitudes be proportional, the greatest and least are greater
than the remaining two.
\end{claim}
\begin{evidence}[Proof of V.25]
\label{ev:V.25}
Let $a : b = c : d$ with $a$ greatest.  By V.14 $b > d$, so
considering $a - c$ and $d - b$ or vice versa, the difference $a - c$
equals $b - d$ in ratio, and rearrangement (with Common Notion 5)
yields $a + d > b + c$.
\dependson{V.25}{V.14}
\dependson{V.25}{V.19}
\dependson{V.25}{cn:5}
\end{evidence}
